\begin{helper}
Fix positive integers $k$ and $D$.  Let $\mathcal H\in H(k,k,D)$ and
let $e\in E(\mathcal H)$.  Then the local value of $(\mathcal H,e)$ is
represented by a bounded local pattern: there is a bounded local pattern
$(\mathcal P,p)$ in the sense of
\noderef{KUniformKSimpleBoundedLocalPattern} such that
\[
b_{L(\mathcal H),kD}(e)\le b_{L(\mathcal P),kD}(p).
\]
\end{helper}
\begin{proof}
Apply \noderef{KUniformKSimpleLocalTruncationPreservesBParameter} to
$\mathcal H$ and $e$.  It gives a local truncation
$(\mathcal H_{\mathrm{loc}},e_{\mathrm{loc}})$ belonging to $H(k,k,D)$,
whose edge labels are exactly $e$ together with the line-graph neighbors of
$e$, whose vertices are exactly the vertices lying in those retained edges,
and whose incidence relation is inherited from $\mathcal H$.  The same node
also gives
\[
b_{L(\mathcal H),kD}(e)
=b_{L(\mathcal H_{\mathrm{loc}}),kD}(e_{\mathrm{loc}}).
\]

It remains to relabel this local truncation into the fixed finite label
sets used by \noderef{KUniformKSimpleBoundedLocalPattern}.  Since
$\mathcal H_{\mathrm{loc}}$ belongs to $H(k,k,D)$,
\noderef{HypergraphClass} gives $k$-uniformity, $k$-simplicity, and maximum
degree at most $D$.  By \noderef{UniformHypergraph}, the distinguished edge
$e_{\mathrm{loc}}$ has exactly $k$ vertices.  Every neighbor of
$e_{\mathrm{loc}}$ meets one of these $k$ vertices, and for each such vertex
\noderef{MaxDegreeAtMost} and \noderef{HypergraphDegree} say that at most
$D$ local edge labels contain it.  Excluding the distinguished edge itself,
there are therefore at most $k(D-1)$ neighbor labels.  Hence the retained
edge-label set has size at most $1+k(D-1)$.

Each retained edge has exactly $k$ vertices by \noderef{UniformHypergraph}.
The retained vertex set is the union of the distinguished edge and the
neighbor edges, so it has size at most
\[
k+k\cdot k(D-1).
\]
Choose injections from the retained edge-label set into
$\{0,\ldots,1+k(D-1)-1\}$ and from the retained vertex set into
$\{0,\ldots,k+k^2(D-1)-1\}$, sending $e_{\mathrm{loc}}$ to the distinguished
label $p$.  Define the bounded pattern $\mathcal P$ by putting, at each used
edge label, the image of the corresponding local edge.  Suppose there are
$m$ unused edge labels.  Then the local truncation has exactly
$1+k(D-1)-m$ retained edge labels, so it has at most
$k+k(k(D-1)-m)$ retained vertices.  The bounded vertex-label set therefore
has at least $km$ labels outside the image of the retained vertex set.  Split
these unused vertex labels into $m$ disjoint $k$-sets and assign one of them
to each unused edge label.  These padding edges are disjoint from the
distinguished edge and from every retained local vertex.  They are
$k$-element edges, have pairwise intersections of size at most $k$, and add
degree one only to newly introduced padding vertices.  Hence the resulting
finite labelled hypergraph still satisfies the three components of
\noderef{HypergraphClass}.

The line graph around the distinguished label is unchanged.  Indeed,
\noderef{LineGraphOfHypergraph} defines adjacency by nonempty hyperedge
intersection, the injections preserve and reflect membership in every used
edge, and the padding edges are disjoint from the distinguished edge, so
they are not neighbors of $p$.  Thus the degree of the distinguished label,
the independent-pair count of its neighborhood from
\noderef{IndependentPairCount}, and the independent-triple count from
\noderef{IndependentTripleCount} are the same for
$(\mathcal H_{\mathrm{loc}},e_{\mathrm{loc}})$ and for
$(\mathcal P,p)$.  By the formula \noderef{LocalBParameter}, their local
$b$-parameters are equal.  Combining this equality with the equality from
the truncation step gives the displayed inequality.
\end{proof}
